\begin{afterword}
\summaryhead

The post takes up Albert's appeal to ``common usage'' in the dispute over whether an unheard falling tree makes a sound. People reach for a dictionary, the author says, because they believe words have meanings that the dictionary records or even sets. But dictionary editors are historians of usage, not lawgivers. The post's example of a dictionary that is properly called ``authoritative'' is the IEEE's, whose terms are negotiated by its members.

The post then grants the strongest objection: shared language is a public good, since we gain from using the same word for the same concept. Yet Albert and Barry each know what the other means, and both have unambiguous words for what happened in the forest. Their appeal to common usage therefore cannot be about communication. If only a name is at stake, they should coin two new words.

Often, though, a category carries hidden inferences or moral stakes, as in whether atheism is a religion or whether blacks and whites are both ``people.'' Then, the post argues, neither common usage nor a dictionary can settle the matter, because definitions used for inference can be wrong, and are wiser or more foolish.

\responsehead

The post's argument is careful. It states the best case for common usage, that shared words are a public good, before it answers it. The answer is that common usage serves communication. When both sides already understand each other, as Albert and Barry do, an appeal to common usage serves nothing, and two new words end the dispute. The post then marks the limit of that remedy: when a category carries inferences or moral stakes, renaming would hide the real question.

The claims about dictionaries hold for the dictionaries the post has in mind. Merriam-Webster says its entries rest on ``usage,'' and the OED calls itself ``a historical dictionary.'' S.~I. Hayakawa made the same argument in 1949, and also documented the belief the post describes, that dictionaries are ``the supreme authority in matters of meaning and usage.'' The post's example of a lawgiving dictionary is the IEEE's; there are older and better-known ones. The French Academy's statutes of 1635 charge it with giving the language fixed rules. And Samuel Johnson, in the preface to his dictionary of 1755, saw a duty to ``correct or proscribe'' improprieties while mocking the hope of fixing the language.

The rule that ``once any empirical proposition is at stake \ldots\ you can no longer appeal to common usage'' is meant, in context, for questions about the things named. Questions about what a speaker or a law meant by a word are empirical too, and there common usage is evidence. In 1893 the Supreme Court decided that tomatoes were vegetables for a tariff act because ``in the common language of the people'' they are, though botanically they are fruit. The court's view of dictionaries, as ``aids to the memory,'' is close to the post's.

In short: a careful account of when an appeal to common usage is idle, whose claims about dictionaries hold for English ones.
\end{afterword}
